% kuleuventheme2poster by Janez Kren, June 2018, janez.kren@kuleuven.be

\documentclass{beamer}
%\usepackage[orientation=portrait,size=a3,scale=1.8,debug]{beamerposter}
\usepackage[orientation=portrait,size=a1,scale=1.8,debug]{beamerposter}
% BEAMERPOSTER OPTIONS:
%	orientation= portrait / landscape
%	size= a0 / a2 / a3 / a4
%	scale = change the size of text (e.g. 1.4 increases all fonts by factor of 1.4)


%\usetheme[sedes,kulak]{kuleuven2poster}
\usetheme[kulak,white]{kuleuven2poster}
% THEME OPTIONS 
%	for logo:   kul (default) / kulak / lrd
%	background colour:   blue (default) / white
%	background sedes logo:   no (default) / sedes
%  e.g. [lrd,white,sedes]


% USE YOUR BACKGROUND IMAGE (optinal):
%\titlegraphic{ \includegraphics[height=0.9\paperheight]{mybackground.jpg} } %or, depending on the size: [width=\paperwidth]
%\def\backgdopacity{0.25}  % background image opacity fraction (0=transparent, 1=full colour)


\usepackage[utf8]{inputenc}
\usepackage{ragged2e}
\usepackage{tikz}
\usepackage{multicol}
\usepackage{algorithm}
\usepackage{algorithmic}
\usepackage{float}


% INFO 
\title{Coloring of planar graphs}
\author{\large Arne Claerhout}
\institute{\large Faculty of Science}
\date{\large Academic year 2025-2026}
 

\newcommand{\kulline}{%
	\par\noindent
	\begin{tikzpicture}[remember picture]
		\draw[kul-blue, opacity=0.3, line width=1] (0,0) -- (\textwidth,0);
	\end{tikzpicture}
	\par
}

\begin{document}
\csname beamer@calculateheadfoot\endcsname %recalculate head and foot dimension

\begin{frame}[t]

%%
%%  TITLE (optional)
%%
\vspace{2.5\baseh} % extra vertical space

\begin{tikzpicture}[remember picture, overlay]
	\node [anchor=north, 
	text width=\textwidth, 
	align=center, 
	inner sep=0pt, 
	outer sep=0, 
	kul-blue, 
	font=\bfseries\fontsize{0.6\baseh pt}{2}\selectfont] 
	at (0.5\textwidth, 0) % Centers the node at the midpoint of the text width
	{ \inserttitle };
\end{tikzpicture}
\vspace{.55\baseh}
%\textcolor{kul-blue}{\bfseries \Huge \inserttitle}



%%
%% BODY
%%
\vspace{.4\baseh}
\begin{columns}[t]
%COLUMN 1
	\begin{column}{.49\textwidth}
	\justifying
	{\bfseries Introduction}
	
	An algorithm for the coloring of planar graphs for multiple coloring variants is proposed. These coloring variants include:
	
	
		\begin{multicols}{2}
			\begin{itemize}
				\item Proper coloring:
				\begin{itemize}
					\item Two adjacent vertices have different colors
				\end{itemize}
			\end{itemize}
			
			\vspace{\fill}\null
			\begin{figure}
				\begin{center}
					\begin{tikzpicture}[
						every node/.style={circle, draw, thick, minimum size=5mm}
						]
						
						\node (v1) at (0, 0) [fill=blue!50]  {};
						\node (v2) at (6, 0) [fill=orange] {};
						\node (v3) at (3, 5.19615) [fill=green]{};
						\node (v4) at (3, 1.73205) [fill=green]{};
						\node (v5) at (3, 3.4641) [fill=red]{};
						
						\draw[thick] (v1) -- (v2);
						\draw[thick] (v1) -- (v3);
						\draw[thick] (v1) -- (v4);
						\draw[thick] (v1) -- (v5);
						
						\draw[thick] (v2) -- (v3);
						\draw[thick] (v2) -- (v4);
						\draw[thick] (v2) -- (v5);
						
						\draw[thick] (v3) -- (v5);
						
						\draw[thick] (v4) -- (v5);
						
					\end{tikzpicture}
				\end{center}
				\caption{A graph colored with the proper coloring.}
				\label{fig:g_proper}
			\end{figure}
			
			
		\end{multicols}
%		\vspace{1cm}
		\kulline
		\begin{multicols}{2}
			\begin{figure}
				\begin{center}
					\begin{tikzpicture}[
						% Define the style for the nodes
						every node/.style={circle, draw, thick, minimum size=5mm}
						]
						
						\node (v1) at (0, 3) [fill=blue!50]  {};
						\node (v2) at (2.853, 0.927) [fill=orange] {};
						\node (v3) at (-2.853, 0.927) [fill=green]{};
						\node (v4) at (1.7634, -2.427) [fill=yellow!40]{};
						\node (v5) at (-1.7634, -2.427) [fill=red]{};
						
						\draw[thick] (v1) -- (v2);
						\draw[thick] (v1) -- (v3);
						\draw[thick] (v2) -- (v4);
						\draw[thick] (v4) -- (v5);
						\draw[thick] (v3) -- (v5);
						
					\end{tikzpicture}
				\end{center}
				\caption{The graph $C_5$ colored with odd coloring.}
				\label{fig:g_odd}
			\end{figure}
			
			\vspace{\fill}\null
			
			\begin{itemize}
				\item Odd coloring
				\begin{itemize}
					\item Proper coloring where every vertex has at least one color appearing an odd amount of times.
				\end{itemize}
			\end{itemize}
			
			\vspace{\fill}\null
			
			
		\end{multicols}
		\vspace{-1cm}
		\kulline
		\begin{multicols}{2}
			\begin{itemize}
				\item Conflict-free coloring
				\begin{itemize}
					\item Every vertex has at least one color occurring only once in its neighborhood.
					\item Four variants: proper/improper and open/closed neighborhoods
				\end{itemize} 
			\end{itemize}
			
			
			\vspace{\fill}\null
			\begin{figure}
				\begin{center}
					\begin{tikzpicture}[
						% Define the style for the nodes
						every node/.style={circle, draw, thick, minimum size=5mm}
						]
						
						% --- VERTICES ---
						% Format: \node (internal_name) at (x,y) [fill=color] {number};
						
						\node (v1) at (0, 0) [fill=blue!50]  {};
						\node (v2) at (6, 0) [fill=orange] {};
						\node (v3) at (3, 5.19615) [fill=green]{};
						\node (v4) at (3, 1.73205) [fill=green]{};
						\node (v5) at (3, 3.4641) [fill=red]{};
						
						% --- EDGES ---
						% In a K4, there are 6 possible edges. 
						% We omit one (e.g., the diagonal between v1 and v4).
						
						\draw[thick] (v1) -- (v2);
						\draw[thick] (v1) -- (v3);
						\draw[thick] (v1) -- (v4);
						\draw[thick] (v1) -- (v5);
						
						\draw[thick] (v2) -- (v3);
						\draw[thick] (v2) -- (v4);
						\draw[thick] (v2) -- (v5);
						
						\draw[thick] (v3) -- (v5);
						
						\draw[thick] (v4) -- (v5);
						
					\end{tikzpicture}
				\end{center}
				\caption{A graph colored for the }
				\label{fig:g_cf}
			\end{figure}
			
			\vspace{\fill}\null
			
		\end{multicols}
		
		\vspace{1cm}
		\kulline
		\begin{multicols}{2}
			
			\begin{figure}
				\begin{center}
					\begin{tikzpicture}[
						every node/.style={circle, draw, thick, minimum size=5mm}
						]
						
						\node (v1) at (0, 3) [fill=blue!50]  {};
						\node (v2) at (2.853, 0.927) [fill=orange] {};
						\node (v3) at (-2.853, 0.927) [fill=green]{};
						\node (v4) at (1.7634, -2.427) [fill=yellow!40]{};
						\node (v5) at (-1.7634, -2.427) [fill=red]{};
						
						\draw[thick] (v1) -- (v2);
						\draw[thick] (v1) -- (v3);
						\draw[thick] (v2) -- (v4);
						\draw[thick] (v4) -- (v5);
						\draw[thick] (v3) -- (v5);
						
					\end{tikzpicture}
				\end{center}
				\caption{A colored graph for the }
				\label{fig:g_um}
			\end{figure}
			
			
			
			\vspace{\fill}\null
			
			\begin{itemize}
				\item Unique-maximum coloring
				\begin{itemize}
					\item The maximum color in the neighborhood of a vertex occurs exactly once.
					\item Has same four variants as conflict-free coloring
				\end{itemize}
			\end{itemize}
			
			
		\end{multicols}
	
	\end{column}


% COLUMNS 2
	\begin{column}{.49\textwidth}
		\justifying
		{\bfseries Chromatic numbers}
		
		\begin{definition}
			Chromatic number $\chi_{K}(\mathcal{P})$ for coloring $K$, is the minimum needed number of colors to color in all graph $G \in \mathcal{P}$ according to $K$.
		\end{definition}
		
		Chromatic numbers for these coloring variants are not always known.
		Conjectures on the chromatic numbers are tested by exhaustively coloring graphs.
		
		\vspace{1cm}
		{\bfseries Algorithm}
		
		Algorithm \ref{alg:search} describes a high-level overview of the designed algorithm.This is a backtracking algorithm using multiple optimizations to perform as efficiently as possible.
		
		\begin{algorithm}[H]
			\caption{\texttt{search\_chromatic\_number}}
			\label{alg:search}
			\textbf{Parameters}: The coloring $K$ and graph $G$
			
			\textbf{Output}: $\chi_{K}(G)$
			\begin{algorithmic}[1]
				\FOR{$max\_c \le MAX$}
				
				\WHILE{not all vertices are colored}
				
				\STATE $v\gets$ \texttt{get\_optimal\_vertex}($G$)
				
				\STATE \texttt{assign\_next\_possible\_color}($v, max\_c$)
				
				\STATE \texttt{update\_neighbors}($G, v$)
				
				\IF{coloring of $G$ is not possible}
					\STATE \texttt{remove\_color}($v$)
				\ENDIF
				
				\ENDWHILE
				
				\ENDFOR
			\end{algorithmic}
		\end{algorithm}
		
		\vspace{0.5cm}
		{\bfseries Results \& Conclusion}
		
		\vspace{0.5cm}
		\begin{block}{Proof of conjectures}
			Conjectures for planar triangulations with number of vertices $n \le 20$ are correct.
		\end{block}
		\begin{itemize}
			\item First open-source algorithm that colors graphs for these variants.
			\item It only performs around $1.4$ times worse than alternatives for the proper coloring. Tools for other colorings do not exist.
		\end{itemize}
		
		
		
		
	\end{column}
\end{columns}

\end{frame}
\end{document}
